\section{All centered trigamma-square moments and a polylogarithmic generator}
\label{ghp:square:sec}

The invariant numerator $(\zeta(2)-H_n^{(2)})^2$ provides an explicit
family beyond raw powers of $H_n$. Define
\begin{equation}\label{ghp:def:square}
 A_n=\zeta(2)-H_n^{(2)}=\psi'(n+1),\qquad
 \mathcal T(s)=\sum_{n=0}^{\infty}\frac{A_n^2}{(n+1/2)^s}.
\end{equation}
The harmonic parity continuation \cite{ghp:report} gives an ordinary finite-sum route to these moments, independent of its broader difference-equation classification. The defining series already converges for $\Re s>-1$. Its centered asymptotic expansion contains only even inverse powers.
Subtracting finitely many terms continues the Dirichlet series normally
to successively larger half-planes, with possible poles only at odd
integers $1-2k$. Hence it is regular at every nonpositive even integer. All its values
there reduce to a rational number plus a rational multiple of
$\zeta(3)$.

\begin{theorem}[Every nonpositive even trigamma-square value]
\label{ghp:thm:square}
For every integer $r\geq0$,
\begin{align}\label{ghp:eq:all-square}
 \mathcal T(-2r)={}&3B_{2r}(1/2)\zeta(3)\notag\\
 &-\frac12\sum_{j=1}^{r}\binom{2r}{2j}
   \frac{2j-1}{2j+1}B_{2r-2j}(1/2)B_{2j-2}.
\end{align}
The Bernoulli number convention is $B_1=-1/2$, and the sum is empty
when $r=0$.
\end{theorem}

For example,
\begin{align}\label{ghp:eq:square-table}
 \mathcal T(0)&=3\zeta(3),\notag\\
 \mathcal T(-2)&=-\frac14\zeta(3)-\frac16,\notag\\
 \mathcal T(-4)&=\frac7{80}\zeta(3)+\frac1{30},\notag\\
 \mathcal T(-6)&=-\frac{31}{448}\zeta(3)+\frac1{672},\notag\\
 \mathcal T(-8)&=\frac{127}{1280}\zeta(3)-\frac{17}{540}.
\end{align}
The first value is the classical Euler sum. The theorem supplies a
single finite formula for all subsequent centered moments.

\subsection{Equivalent identities of absolutely convergent ordinary sums}

Let
\begin{equation}\label{ghp:def:ck}
 c_k=\sum_{i=0}^{k-1}B_{2i}(1/2)B_{2k-2-2i}(1/2),\qquad k\geq1.
\end{equation}
Thus $c_1=1$, $c_2=-1/6$, $c_3=47/720$, and
$c_4=-257/5040$. The shifted Stirling expansion gives $A_n\sim\sum_{i\ge0} B_{2i}(1/2)(n+1/2)^{-2i-1}$. Squaring this expansion gives
$A_n^2\sim\sum_{k\geq1}c_k(n+1/2)^{-2k}$.
For every $r\geq0$ the exact ordinary convergent identity is
\begin{equation}\label{ghp:eq:convergent-square}
 \sum_{n=0}^{\infty}\left\{(n+1/2)^{2r}\psi'(n+1)^2
       -\sum_{k=1}^r c_k(n+1/2)^{2r-2k}\right\}
       =\mathcal T(-2r),
\end{equation}
where the right side is explicitly \eqref{ghp:eq:all-square}. The
summands after subtraction are $O(n^{-2})$.

For instance,
\[
 \sum_{n=0}^{\infty}\left((n+1/2)^2\psi'(n+1)^2-1\right)
       =-\frac14\zeta(3)-\frac16,
\]
and
\[
 \sum_{n=0}^{\infty}\left((n+1/2)^4\psi'(n+1)^2
                         -(n+1/2)^2+\frac16\right)
       =\frac7{80}\zeta(3)+\frac1{30}.
\]
These statements need no prescription for summing a divergent series.
To prove \eqref{ghp:eq:convergent-square}, subtract the displayed
asymptotic terms in the half-plane of convergence and continue. Their
Dirichlet sums at $s=-2r$ are
$c_k\zeta(2k-2r,1/2)=0$, since
$\zeta(-2m,1/2)=0$ for $m\geq0$.

\subsection{Proof of the Bernoulli formula by finite summation}

For a sequence admitting an expansion in powers of $N$ and $\log N$,
write $\operatorname{CT}_N$ for its coefficient of $N^0(\log N)^0$.
Only finite exact identities followed by known asymptotic expansions
will be used here. Put
\begin{equation}\label{ghp:def:S}
 S_r(k)=\sum_{n=0}^{k-1}(n+1/2)^{2r}
  =\frac{B_{2r+1}(k+1/2)}{2r+1}
  =\sum_{j=0}^r b_{r,j}k^{2j+1},
\end{equation}
where
\begin{equation}\label{ghp:def:b}
 b_{r,j}=\frac{\binom{2r}{2j}}{2j+1}B_{2r-2j}(1/2).
\end{equation}
The Bernoulli polynomial $B_{2r+1}(1/2)$ is zero, so this is an odd
polynomial in $k$ and has zero constant term.

\begin{proof}[Proof of Theorem~\ref{ghp:thm:square}]
First, \eqref{ghp:eq:convergent-square} and \eqref{ghp:def:S} show that
\begin{equation}\label{ghp:eq:ct-reduction}
 \mathcal T(-2r)=\operatorname{CT}_N
       \sum_{n=0}^{N-1}(n+1/2)^{2r}A_n^2.
\end{equation}
Indeed every subtracted even power has partial sum with zero constant
term in $N$, and the convergent remainder tends to the stated value.

The identity $A_{k-1}=A_k+k^{-2}$ gives finite summation by parts:
\begin{align}\label{ghp:eq:partial-sum}
 \sum_{n=0}^{N-1}(n+1/2)^{2r}A_n^2
 ={}&\sum_{k=1}^N S_r(k)
          \left(\frac{2A_k}{k^2}+\frac1{k^4}\right)
       +S_r(N)A_N^2.
\end{align}
The coefficient $b_{r,0}$ of $k$ contributes
\[
 \operatorname{CT}_N
 \left(2\sum_{k=1}^N\frac{A_k}{k}+H_N^{(3)}+NA_N^2\right)
 =3\zeta(3).
\]
Here $\sum_{k\geq1}A_k/k=\zeta(2,1)=\zeta(3)$ is the classical
Euler identity and $NA_N^2\to0$. For completeness, the Euler identity
follows from
\[
 \zeta(2,1)=\frac12\int_0^1\frac{\log^2(1-t)}t\,dt
 =\frac12\int_0^\infty\frac{u^2}{e^u-1}\,du=\zeta(3).
\]

For $j\geq1$, it remains to compute the constant of
\begin{equation}\label{ghp:def:Ej}
 E_j(N)=2\sum_{k=1}^Nk^{2j-1}A_k
        +\sum_{k=1}^Nk^{2j-3}+N^{2j+1}A_N^2.
\end{equation}
Let
\[
 F_j(N)=\sum_{k=1}^Nk^{2j-1}
       =\frac{B_{2j}(N+1)-B_{2j}}{2j}.
\]
Changing the order in a finite sum gives
\begin{equation}\label{ghp:eq:finite-change-order}
 \sum_{k=1}^Nk^{2j-1}A_k
  =A_NF_j(N)+\sum_{\ell=1}^N\frac{F_j(\ell-1)}{\ell^2}.
\end{equation}
For $j\geq2$, the second summand is a sum of polynomial powers of
$\ell$ and has zero constant term as a polynomial in $N$;
$B_{2j-1}=0$ removes a possible $1/\ell$ term. The other power sum in
\eqref{ghp:def:Ej} also has zero constant. For $j=1$ the second sum in
\eqref{ghp:eq:finite-change-order} equals $(N-H_N)/2$; its harmonic
term cancels exactly against $\sum k^{-1}$ in \eqref{ghp:def:Ej}.
In all cases,
\begin{equation}\label{ghp:eq:Ej-CT}
 \operatorname{CT}_NE_j(N)
   =2\operatorname{CT}_N(A_NF_j(N))
          +\operatorname{CT}_N(N^{2j+1}A_N^2).
\end{equation}

Now use the uncentered tail expansion
\[
 A_N=\frac1N-\frac1{2N^2}
                  +\sum_{k\geq1}\frac{B_{2k}}{N^{2k+1}}.
\]
The finite coefficient calculation gives
\begin{equation}\label{ghp:eq:two-constants}
 \operatorname{CT}_N(A_NF_j(N))
      =\frac{3-2j}{4}B_{2j-2},\qquad
 \operatorname{CT}_N(N^{2j+1}A_N^2)=-B_{2j-2}.
\end{equation}
To check the first identity, the coefficient of $N^\ell$ in $F_j$
is $\binom{2j}{\ell}B_{2j-\ell}(1)/(2j)$. Only $\ell=2$ and
$\ell=2j-1$ can pair with nonzero inverse coefficients of $A_N$
when $j\geq2$; their contributions are
$-(2j-1)B_{2j-2}/4$ and $B_{2j-2}/2$. The case $j=1$ is the direct
product with $F_1(N)=(N^2+N)/2$ and gives $1/4$.
For the second identity, an odd inverse exponent in $A_N^2$ arises
only by multiplying $-1/(2N^2)$ by the appropriate odd inverse power;
this gives $-B_{2j-2}$, also at $j=1$.

Combining \eqref{ghp:eq:Ej-CT} and \eqref{ghp:eq:two-constants} yields
\[
 \operatorname{CT}_NE_j(N)=-\frac{2j-1}{2}B_{2j-2}.
\]
Multiplication by $b_{r,j}$, followed by the $j=0$ contribution,
proves \eqref{ghp:eq:all-square}.
\end{proof}

\subsection{A convergent generator involving classical polylogarithms}

The complete sequence has an analytic exponential generating function.
Let
\[
 \mathcal G(t)=\sum_{r=0}^{\infty}
               \mathcal T(-2r)\frac{t^{2r}}{(2r)!}.
\]

\begin{theorem}[Polylogarithmic moment generator]\label{ghp:thm:generator}
The series $\mathcal G$ converges for $|t|<2\pi$ and equals
\begin{equation}\label{ghp:eq:integral-generator}
 \mathcal G(t)=\frac{t}{2\sinh(t/2)}
   \left[3\zeta(3)-\frac1{4t}
       \int_0^t(t-u)u^2\coth(u/2)\,du\right].
\end{equation}
The singularities at zero in this expression are removable. For
real $0<t<2\pi$, an equivalent formula in classical polylogarithms is
\begin{align}\label{ghp:eq:polylog-generator}
 \mathcal G(t)=\frac{t}{2\sinh(t/2)}\bigg[
 2\zeta(3)-\frac{t^3}{48}
 -\frac t2\operatorname{Li}_2(e^{-t})
 -2\operatorname{Li}_3(e^{-t})
 +\frac{3\bigl(\zeta(4)-\operatorname{Li}_4(e^{-t})\bigr)}t
 \bigg].
\end{align}
Its continuation specified by \eqref{ghp:eq:integral-generator} is
holomorphic and even throughout $|t|<2\pi$.
\end{theorem}

\begin{proof}
The Bernoulli generating functions give
\[
 \frac{t}{2\sinh(t/2)}
     =\sum_{r\geq0}B_{2r}(1/2)\frac{t^{2r}}{(2r)!},
 \qquad
 \frac u2\coth(u/2)
     =\sum_{k\geq0}B_{2k}\frac{u^{2k}}{(2k)!}.
\]
Both converge in the disk of radius $2\pi$. Integrating the second
series twice gives
\[
 \frac1{2t}\int_0^t(t-u)u^2\coth(u/2)\,du
  =\sum_{j\geq1}\frac{2j-1}{2j+1}B_{2j-2}
                  \frac{t^{2j}}{(2j)!}.
\]
Thus \eqref{ghp:eq:integral-generator} is exactly the binomial
convolution in \eqref{ghp:eq:all-square}. This proves convergence,
analyticity and the integral formula. In particular all apparent odd
or logarithmic terms in any branch representation must cancel.

For positive $t$, expand $(e^u-1)^{-1}$ in its convergent positive
exponential series. Direct integration, justified for the integrable
factor $(t-u)u^2$, gives
\begin{align*}
 \int_0^t\frac{(t-u)u^2}{e^u-1}\,du
  ={}&2t\zeta(3)-6\zeta(4)
      +t^2\operatorname{Li}_2(e^{-t})\\
     &+4t\operatorname{Li}_3(e^{-t})
      +6\operatorname{Li}_4(e^{-t}).
\end{align*}
Finally $\coth(u/2)=1+2/(e^u-1)$ and
$\int_0^t(t-u)u^2\,du=t^4/12$ give
\eqref{ghp:eq:polylog-generator}.
\end{proof}


The finite-sum proof uses no transcendence assumption and no formal
difference Galois theorem. The subtractions in the ordinary sums are part
of each identity; without them the higher moments would diverge. The
broader reflection-invariant numerator classification, resolvent transforms
and independently marked higher-tail moments retain separate audit scopes.
